\chapter{Surface Waves}\label{exten}


\section{The elastodynamic Jost Solutions in a stratified 
half space}\label{subsecJost}

Solving the elastodynamic wave equations is greatly simplified by 
exploiting their {\em linearity\/}. The solution of a linear problem, may 
expressed in terms of the {\em basis states\/} of the solution space. In 
seismology, the solution space basis is referred to as the {\em 
fundamental solution\/}. When  modeling wave propagation in a plane stratified medium, the 
{\em plane wave basis\/} is pre-eminent, because plane 
stratified media are transversely invariant, and plane waves are the 
eigen vectors of the translation operator (an eigen basis is always the 
preferred basis, assuming that it is available). Each of these plane waves 
is a solution to the wave equation in its own right.

For  fixed values of $(\omega, k_y)$, the  elastodynamic basis states 
may be classified on the basis of $3$ separate 
physical properties:
\begin{itemize}
\item the scattering characteristics of an earth region can be defined by the outcome of a {\em 
dual pair of scattering experiments}, initiated by downward incident 
energy from above, and upward incident energy from below, labeled as experiment ``$u$'' and ``$v$''
respectively. We shall consider these scattering experiments in detail later, but for now we
merely point out that they are the elastodynamic generalization of a related pair of
scattering experiments commonly employed in quantum scattering. In the solution of the
quantum problem, the outcomes are referred to as the {\em Jost solutions}, hence we
adopt the appellation  {\em  elastodynamic Jost solution} to label these outcomes in the
elastodynamic problem. When the reference medium is a homogeneous space, 
rather than a stratified half space,  the ``$u$'' solution is exclusively 
comprised of downgoing plane waves, while the ``$v$'' solution is exclusively 
comprised of upgoing plane waves. But when the reference medium is 
stratified,  ``$u$'' and ``$v$'' have both up and downgoing components.
\item each such elastodynamic Jost solution is further characterized by its
incident polarization  i.e. either  {\em P},  {\em SV} or {\em SH-}wavemotion
\item  at any depth $z$  the transverse (i.e. $x$) dependence  of the wavefield is expanded 
in a Fourier basis of plane 
waves.  The incident field is further classified on the basis of the incident 
horizontal wavenumber $k_{x}$
\end{itemize}
We introduce {\em multi-indices} $\nu=(J,c,k_x)\in \{ \{u,v  \} \times \{ 
P, S, H \} 
\times (-\infty, \infty)   \}$ to summarize these three 
physical properties, i.e. the first index $J$ takes the value $J=u$ for the 
``$u$'' Jost solution and $J=v$ for the ``$v$'' Jost solution. The second 
index $c$ , takes values $P$, $S$ and $H$, denoting incident {\em P} 
{\em SV} and {\em SH-}polarizations respectively.  The final index $k_x$ of the 
3-tuple multi-index is the
horizontal wavenumber component, so that there are $6$ different basis states 
for each value of $k_x$.

 The solution basis states (or partial wave 
states) at ${\bmi x}=(z,x)$ and fixed $(\omega, k_y)$ may be written in the Dirac notation as
\[
\left|\nu, {\bmi x}  \right\rangle, \quad \nu=(J,c,k_x)\in \{ \{u,v  \} \times \{ 
P, S, H \} 
\times (-\infty, \infty)   \}
\]
If we introduce the synthesis operator
\[
\sum\nolimits_\nu\!\!\!\!\!\!\!\!\! \int\,\,\, \equiv {\frac{1}{2 \pi }}\int_{- \infty }^{ \infty }\, {\rm d}{k}_{x}\, 
\sum\limits_{ J  \in   \left\{{u,v}\right\}} 
\sum\limits_{ c  \in   \left\{{P,S,H}\right\}}
\]
then, at fixed $(\omega, k_y)$, we synthesize an 
arbitrary displacement/ traction field 
\begin{eqnarray} \label{bdefSurf}
\lefteqn{{\bmi b}({\bmi x};k_y,\omega)=}\\
&&\!\!\!\!\!\!
\left[u_{z}({\bmi x};k_y,\omega),\ u_{x}({\bmi x};k_y,\omega),\ u_{y}({\bmi x};k_y,\omega), 
\ \tau_{zz}({\bmi x};k_y,\omega),\  \tau_{xz}({\bmi x};k_y,\omega),\ 
\tau_{yz}({\bmi x};k_y,\omega) \right]^{T}\nonumber
\end{eqnarray}
in a stratified reference medium, from the partial 
wave states according to
\begin{equation} \label{bsynthesis}
{\bmi b}({\bmi x};k_y,\omega)=
\sum\nolimits_\nu\!\!\!\!\!\!\!\!\! \int \,\,\, \left|\nu, {\bmi x}  
\right\rangle.
\end{equation}
 
The {\em partial wave synthesis\/} (\ref{bsynthesis}) is a rather concise statement of a 
complex physical process, but has little practical utility unless we know 
how the partial wave states $\left|\nu, {\bmi x}  
\right\rangle$ evolve over space. We have already met this problem under 
various guises in the preceding chapters. We recall the spectral 
representation (\ref{Ddef}), where we constructed the fixed 
$(k_x,k_y,\omega)$ displacement/ traction field
${\bmi b}(k_x,k_y,\omega)$ in free space, using the eigen states of the system 
matrix (\ref{Akdef}). In this section, we will exploit the plane 
stratified nature of the reference medium, so that the $k_x$ label
of a partial wave state  is conserved as the state evolves over depth.
Consequently, the spectral 
representation (\ref{Ddef}) may be used over the depth interval $z_{l-1} \le z \le 
z_{l}$  of plane reference layer$-l$ for each $1 \le l \le L$, but we 
must allow for coupling between the different $c\in \{P, S, H\}$ labels 
as we cross the plane interface between adjacent layers. Of course, this 
problem has been studied in great detail in numerous texts on 
theoretical seismology.  



Due to the plane stratified nature of the reference medium, ${\bmi 
b}({\bmi x};k_y,\omega)$ may be constructed over the depth interval $z_{l-1} \le z \le 
z_{l}$  of plane reference layer$-l$ for each $1 \le l \le L$,  according to
\begin{eqnarray} \label{brepj}
{\bmi b}_{l}({\bmi x};k_y,\omega)&=& \sum\nolimits_{\nu_{l}}\!\!\!\!\!\!\!\!\!\!\! 
\int \,\,\,\,\, \left|\nu_{l}, {\bmi x}  
\right\rangle\\
&=& \frac{1}{2 \pi} \int_{-\infty}^{\infty} {\rm d}k \, 
{\bmi \Theta}_{l}({\bmi x};k_{x},k_{y})
\left[\begin{array}{c}
{\bmi V}_{\downarrow 
,P_{l} }({z}_{l-1};k_{x},k_{y})\\
{\bmi V}_{\downarrow ,S_{l} }({z}_{l-1};k_{x},k_{y})\\
{\bmi V}_{\downarrow ,H_{l} }({z}_{l-1};k_{x},k_{y})\\
{\bmi V}_{\uparrow , P_{l} }({z}_{l};k_{x},k_{y})\\
{\bmi V}_{\uparrow , S_{l} }({z}_{l};k_{x},k_{y})\\
{\bmi V}_{\uparrow , S_{l} }({z}_{l};k_{x},k_{y})
\end{array}\right], \quad z_{l-1}\le z \le z_{l}\nonumber
\end{eqnarray}
where:
\begin{itemize}
\item the partial wave amplitudes 
\[{\bmi V}_{\downarrow ,P_{l} 
}({z}_{l-1}; k_{x},k_{y}),\quad 
{\bmi V}_{\downarrow ,S_{l} }({z}_{l-1};k_{x},k_{y})\quad \mbox{and}
\quad{\bmi V}_{\downarrow ,H_{l} }({z}_{l-1};k_{x},k_{y})
\]
 are the amplitudes 
of the downgoing  $P$, $SV$ and $SH$  partial waves  respectively, phased relative to 
the top of layer-$l$
\item the partial wave amplitudes 
\[{\bmi V}_{\uparrow ,P_{l} }({z}_{l};k_{x},k_{y}),\quad
{\bmi V}_{\uparrow ,S_{l} }({z}_{l};k_{x},k_{y})\quad \mbox{ and} 
\quad{\bmi V}_{\uparrow ,H_{l} }({z}_{l};k_{x},k_{y})
\]
 are the amplitudes of the upgoing $P$, $SV$ and $SH$ partial waves  respectively, phased relative to 
the bottom of layer-$l$
\item the $6 \times 6$ array
\begin{eqnarray} \label{Thetaldef}
\lefteqn{{\bmi \Theta}_{l}({\bmi x};k_{x},k_{y})=
{\rm e}^{i{ k}_{x} x} {\bmi D}_{z,l}(k_{x},k_{y})}\\
&&\!\!\!\!\!\!
\times
\mbox{diag}\left[\begin{array}{cccccc}
{\rm e}^{i{ k}_{z,P_l} (z-z_{l-1})}&
{\rm e}^{i{ k}_{z,S_l} (z-z_{l-1})}
&{\rm e}^{i{ k}_{z,H_l} (z-z_{l-1})}
&{\rm e}^{i{ k}_{z,P_l} (z_{l}-z)}
&{\rm e}^{i{ k}_{z,S_l} (z_{l}-z)}&
{\rm e}^{i{ k}_{z,H_l} (z_{l}-z)}
\end{array}\right]\nonumber
\end{eqnarray} 
{\em synthesizes the  displacement-traction field\/} at fixed $(k_x,k_y, 
\omega)$ from the partial wave amplitudes

\item  ${\bmi D}_{z}(k_{x},k_{y})$  was given by (\ref{Ddef})
\item ${ k}_{c_l,z}$ was given by (\ref{kzc1Def})
\end{itemize}

Using (\ref{Thetaldef}), we now construct a $6 $ element basis 
${\bmi B} ({\bmi x};k_{x},k_{y})$ of solutions like (\ref{brepj}) 
for fixed $(k_{x},k_{y},\omega)$, using the partial wave amplitudes 
${\bmi V}_{\downarrow,\uparrow ,c_{l} 
}$ within each layer. 
This fundamental solution ${\bmi B} ({\bmi x};k_{x},k_{y})$ may be 
written  as
\begin{eqnarray} \label{fundamental}
\lefteqn{{\bmi B} ({\bmi x};k_{x},k_{y}) = 
\left[\begin{array}{cc} {\bmi  B}_{u} ({\bmi x};k_{x},k_{y}) & {\bmi  
B}_{v} ({\bmi x};k_{x},k_{y})\end{array}\right]}\nonumber\\
&&= {\bmi \Theta}_{l}({\bmi x};k_{x},k_{y})
\left[\begin{array}{cc} {\bmi  V}_{u,l}(k_{x},k_{y}) & {\bmi  
V}_{v,l}(k_{x},k_{y})\end{array}\right], \quad     z_{l-1}\le z \le z_{l}, \quad 
 1 \le l \le L\nonumber\\
 &&= {\bmi \Theta}({\bmi x};k_{x},k_{y})\left[\begin{array}{cc} 
 {\bmi  V}_{u}(k_{x},k_{y}) & 
 {\bmi  V}_{v}(k_{x},k_{y})\end{array}\right]
\end{eqnarray} 
where:
\begin{itemize}
\item 
the  $6L\times 6 L$ block diagonal displacement-traction synthesis array 
${\bmi \Theta}({\bmi x};k_{x},k_{y})$ at fixed $(k_x,k_y,\omega)$,
 is defined in terms of its layer contributions 
${\bmi \Theta}_{l}({\bmi x};k_{x},k_{y}), 1\le l \le L$, from 
(\ref{Thetaldef}) as
\[
{\bmi \Theta}({\bmi x};k_{x},k_{y})\equiv
\mbox{diag}\left[\begin{array}{ccc}{\bmi \Theta}_{1}({\bmi x};k_{x},k_{y})&
\cdots& {\bmi \Theta}_{L}({\bmi x};k_{x},k_{y})  \end{array}  \right]
\]
\item the $(6  L)\times (3 )$ arrays 
\begin{equation} \label{Vvw} 
\!\!\!\!\!\!
{\bmi {V}}_{u}(k_{x},k_{y})=
\left(\begin{array}{c} 
{{\bmi V}}_{ u,1}(k_{x},k_{y})\\
\vdots\\ {{\bmi V}}_{ u,l}(k_{x},k_{y})\\
\vdots\\  {{\bmi V}}_{ u,L}(k_{x},k_{y})
\end{array}\right) \ \mbox{and} \
{\bmi {V}}_{v}(k_{x},k_{y})=
\left(\begin{array}{c} 
{{\bmi V}}_{ v,1}(k_{x},k_{y})\\
\vdots\\ {{\bmi V}}_{ v,l}(k_{x},k_{y})\\
\vdots\\  {{\bmi V}}_{ v,L}(k_{x},k_{y})
\end{array}\right)
\end{equation}
are to be constructed from the matrix factorizations of \S\ref{RTfactor}, after we 
have chosen the appropriate source radiation
\end{itemize}
We propose that:
\begin{itemize}
\item ${\bmi  V}_{u}(k_{x},k_{y})$,  
is initiated by a {\em free surface source layer\/}
with source radiation  \newline
${\bmi \Sigma}_{u}(k_{x},k_{y}) = [{\bmi 
\Sigma}_{u,\uparrow}(k_{x},k_{y}),{\bmi 
\Sigma}_{u,\downarrow}(k_{x},k_{y})]^{T} =
[{\bf 0}_{3},{\bmi I}_{3}]^{T}$ 
radiating $P$  waves, $SV$ waves and $SH$ waves down into the stratified half space. 
Each such partial wave component of ${\bmi  V}_{u}(k_{x},k_{y})$ satisfies the 
radiation condition in the lower half space. 
\item ${\bmi  V}_{v}(k_{x},k_{y})$, has its source in the lower half space
with source radiation \newline
${\bmi \Sigma}_{v}(k_{x},k_{y})=
[{\bmi \Sigma}_{v,\uparrow}(k_{x},k_{y}),
{\bmi \Sigma}_{v,\uparrow}(k_{x},k_{y})]^{T}=[{\bmi I}_{3},
{\bf 0}_{3}]^{T}$ radiating $P$  waves, $SV$ waves and $SH$  waves from the lower half space. Each such partial wave component of 
${\bmi  V}_{v}(k_{x},k_{y})$ satisfies the 
{\em traction free condition\/} at the free surface.
\end{itemize}
In the downward reflection, and upward transmission problems of this 
thesis, 
we are interested in modeling wave propagation on the interval $z_0\le 
\zeta \le z_L$, so that the source depths $z_S$ for 
${\bmi  V}_{u}$ and ${\bmi  V}_{v}$ are chosen outside this interval. 
Otherwise, if $z_{0}\le 
z_S \le z_{L}$, then ${\bmi  V}_{u}$ and ${\bmi  V}_{v}$ will suffer jump 
discontinuities across  $z_S$ and this will complicate the analysis 
considerably.  Consequently, the source fields for ${\bmi  V}_{u}$ and ${\bmi  V}_{v}$
irradiate the layer stack from above, and below respectively.
The fields ${\bmi  B}_{u}({\bmi x};k_{x},k_{y})$ and 
${\bmi  B}_{v}({\bmi x};k_{x},k_{y})$ in (\ref{fundamental}), play a 
role analogous to the matrix {\em Jost solutions\/} in mathematically rigorous 
formulations of quantum scattering [\cite{Taylor}, \cite{Newton}] and also 
the theory of solitons \cite{Lamb}. The precise definition of the source 
radiation ${\bmi \Sigma}_{u}(k_{x},k_{y})$ and 
${\bmi \Sigma}_{v}(k_{x},k_{y})$, initiating the two Jost solutions 
will not affect the final form of the equations, 
but may simplify the development.


By setting source depth $z_s=f+$ (i.e. just below the free surface) and  
source radiation ${\bmi \Sigma}(k_{x},k_{y})={\bmi \Sigma}_{u}
(k_{x},k_{y})=[{\bmi 
\Sigma}_{u,\uparrow}(k_{x},k_{y}),{\bmi \Sigma}_{u,\downarrow}
(k_{x},k_{y})]^{T}=
[{\bf 0}_{3 },{\bmi I}_{3 }]^{T}$  in (\ref{incdown}), 
the depth evolution of ${\bmi {V}}_{u}(k_{x},k_{y})$, may be written in terms of its reference layer 
contributions as
\begin{eqnarray} \label{Vudef} 
{{\bmi V}}_{u,l}(k_{x},k_{y})&=&
\left[\begin{array}{c}
{\bmi I}_{3 }\\ 
{\bmi S}{}_{\uparrow \downarrow }^{ [l , \infty)}
 \ {{\bmi E}}_{l}\end{array}\right] 
\ {\left[{ \ {\bmi  I}_{3 }-
{\bmi S}{}_{\downarrow \uparrow }^{ (f , l-1]} \ {\bmi S}{}_{ \uparrow 
\downarrow }^{ (l-1 , \infty)} \ }\right]}^{  - 1}
{\bmi S}{}_{\downarrow \downarrow }^{ (f , l-1]} \nonumber\\
&& \qquad \qquad \times  {\left[{ {\bmi  I}_{3 } 
- {\bmi S}{}_{ \downarrow \uparrow }^{ f} \ {\bmi S}{}_{\uparrow \downarrow 
}^{(f, \infty)}  }\right]}^{ - 1}(k_{x},k_{y}),\qquad 1 \le l \le L
\end{eqnarray}
and ${\bmi S}{}_{\downarrow \uparrow }^{ (f , l-1]}(k_{x},k_{y})$ is the net upward 
incident reflection operator for the half open depth interval
$0< z \le l-1$ from the top of layer-$l$, but with 
the free surface reflection operator 
${\bmi S}{}_{ \downarrow \uparrow }^{ f}(k_{x},k_{y})$ set to zero.



 
The interface variables ${\bmi V}_{ u,l}(k_{x},k_{y})$ of
(\ref{Vudef}) characterize the solution field in each of $L$ uniform 
layers by the amplitudes of $ \downarrow \!  \!-going\ $
 {\em P, SV\/} and {\em SH\/} plane waves phased relative to  the top of layer-$l$ (ie. 
\( z={z}_{l-1}^{+}\) ), and the amplitudes of
\(  \uparrow \!  
\!-going\ \) 
{\em P, SV\/} and {\em SH\/} waves phased 
relative to the bottom of layer-$l$ (i.e.
\(z={z}_{l}^{-}\) ) for each   $1 \le    l \le   L$
 of the reference stratification .


Using ${\bmi V}_{ u,l}(k_{x},k_{y})$ from (\ref{Vudef}), the {\em 
Jost solution\/} ${\bmi B}_{u}$ is written as
\begin{eqnarray}\label{Budef}
{\bmi B}_{u}({\bmi x};k_{x},k_{y})
&\equiv&{\rm e}^{i{ k}_{x} x}\,
\left[\begin{array}{ccc} (u,P,k_x,k_{y},z)&(u,S,k_x,k_{y},z)
&(u,H,k_x,k_{y},z)\end{array} \right]\nonumber\\
&\equiv&{\bmi \Theta}_{l}({\bmi x};k_{x},k_{y}){{\bmi 
V}}_{u,l}(k_{x},k_{y}),
\ \  z_{l-1} \le z \le{z}_{l} \ \mbox{for each}  \ 1 \le    l \le   L
\end{eqnarray}
where ${\bmi \Theta}_{l}({\bmi x};k_{x},k_{y})$ was defined in 
(\ref{Thetaldef}).
	
Similarly, from (\ref{incup}),  with the source depth $z_s=z_{L}+$ (i.e. 
at the top of the lower half-space, {\em see\/} Fig.\ref{BasicStack}), the  source radiation 
${\bmi \Sigma}(k_{x},k_{y})={\bmi \Sigma}_{u}(k_{x},k_{y})=[{\bmi 
\Sigma}_{u,\uparrow}(k_{x},k_{y}),{\bmi \Sigma}_{u,\downarrow}(k_{x},k_{y})]^{T}=
[{\bmi I}_{3 },{\bf 0}_{3 }]^{T}$ and 
${\bmi S}{}_{\uparrow\downarrow}^{(z_s, \infty)}(k_{x},k_{y})={\bf 0}_{3 }$ (because we now 
assume for simplicity, that the medium is homogeneous below the source),   
the depth evolution of ${\bmi V}_{v}$, may be written in terms of its reference layer 
contributions as
\begin{equation} \label{Vvdef} 
{{\bmi V}}_{ v,l}(k_{x},k_{y})=
\left[{\begin{array}{c} 
{\bmi S}{}_{\downarrow \uparrow}^{[f, l-1]}\ {{\bmi E}}_{l} 
\\{\bmi I}_{3 }\end{array}}\right]
 {\left[{  {\bmi  I}_{3 }- 
 {\bmi S}{}_{ \uparrow \downarrow }^{[l, \infty)} 
 {\bmi S}{}_{\downarrow \uparrow }^{[f, l)} }\right]}^{-1}
 {\bmi S}{}_{\uparrow \uparrow }^{ [l, L]}(k_{x},k_{y}), \qquad 1 \le l \le L .
\end{equation}

Using ${\bmi V}_{ v,l}(k_{x},k_{y})$ from (\ref{Vvdef}), the {\em 
Jost solution\/} ${\bmi B}_{v}$ is written as
\begin{eqnarray}\label{Bvdef}
{\bmi B}_{v}({\bmi x};k_{x},k_{y})
&\equiv&{\rm e}^{i{ k}_{x} x}\,\left[\begin{array}{ccc} 
(v,P,k_x,k_{y},z)&(v,S,k_x,k_{y},z)
&(v,H,k_x,k_{y},z)\end{array} \right]\nonumber\\
&\equiv&{\bmi \Theta}_{l}({\bmi x};k_{x},k_{y}){{\bmi V}}_{v,l}(k_{x},k_{y})
\  z_{l-1} \le z \le{z}_{l}, \ \ \mbox{for each} \  1 \le    l \le   L
\end{eqnarray}
where ${\bmi \Theta}_{l}({\bmi x};k_{x},k_{y})$ was defined in 
(\ref{Thetaldef}).
\subsection{Elementary properties of the Jost solutions}

We now introduce  the {\em skew symmetric bilinear functional\/} 
$\langle \cdot , \cdot \rangle$, given by
\begin{equation} \label{functional}
\langle {\bmi  u}_{1} , {\bmi  u}_{2} \rangle \stackrel{\mbox{\tiny def}}{=}
({\bmi  u}_{1})^{ T} {\bmi  J}_6 {\bmi  u}_{2}
\end{equation}
for a pair of $6-$vector fields ${\bmi  u}_{1}$ and ${\bmi  
u}_{2}$.
If instead, we have a $6\times l$ array
${\bmi V}=({\bmi u}_{1}  ,\cdots,  {\bmi u}_{l})$ and a $6\times m$ array
${\bmi W}=({\bmi v}_{1}  ,\cdots,  {\bmi v}_{m})$
then
\begin{equation} \label{functionalnm}
\langle {\bmi V} , {\bmi W}
\rangle= 
\left(
\begin{array}{ccc}
\langle {\bmi u}_{1},  {\bmi v}_{1} \rangle&\cdots&\langle {\bmi u}_{1},  {\bmi
v}_{m} \rangle\\
\cdots&\langle {\bmi u}_{i},  {\bmi v}_{l} \rangle&\cdots\\
\langle {\bmi u}_{l},  {\bmi v}_{1} \rangle&\cdots&\langle {\bmi u}_{l},  {\bmi
v}_{m} \rangle
\end{array}
\right)
\end{equation}
is an $l \times m$ array, where $\langle {\bmi u}_{i},  {\bmi v}_{l} 
\rangle$ is defined by (\ref{functional}). 

In the next section, we will evaluate a number of functionals similar to 
the prototype
\begin{equation}\label{BuBv}
\langle {\bmi B}_{u}({\bmi x};-k_{x},-k_{y}) , 
{\bmi B}_{v}({\bmi x};k_{x},k_{y})\rangle
\end{equation}
 where ${\bmi B}_{u}({\bmi x};-k_{x},-k_{y}) $ and 
$ {\bmi B}_{v}({\bmi x};k_{x},k_{y})$ are the {\em Jost solutions\/} of 
(\ref{Budef}) and  (\ref{Bvdef})  respectively.
The key to the efficient evaluation of (\ref{BuBv})
is to observe that it is depth invariant over the interval $f+\le z \le 
L+$ because ${\bmi B}_{u}({\bmi x};-k_{x},-k_{y})$ and 
${\bmi B}_{v}({\bmi x};k_{x},k_{y})$ are {\em 
homogeneous\/} solutions of (\ref{Leqn})
\begin{equation} \label{LeqnSurf}
{\bmi L}(z,x;{\partial }_{z},{\partial }_{x},{k}_{y})
{\bmi B}_{v}({\bmi x};k_{x},k_{y}) = 
{\bmi L}(z,x;{\partial }_{z},{\partial }_{x},-{k}_{y})
{\bmi B}_{u}({\bmi x};-k_{x},-k_{y}) ={\bf 0}
\end{equation}
where 
\[
{\bmi L}(z,x;{\partial }_{z},{\partial }_{x},{k}_{y})\equiv
\left[\partial_{z}-\rf{\bmi A}(z,x;{\partial }_{x},{k}_{y})\right]
\]
was given earlier in (\ref{Lref}).
Using (\ref{LeqnSurf}) we have
\begin{eqnarray} \label{invariant}
\lefteqn{
{\partial }_{z}\left\langle{\bmi B}_{u}({\bmi x};-k_{x},-k_{y}),
{\bmi B}_{v}({\bmi x};k_{x},k_{y})\right\rangle}\nonumber\\
&&\!\!\!\!\!\!
=
\left\langle{{ \partial 
}_{z}{\bmi B}_{u}({\bmi x};-k_{x},-k_{y}),
{\bmi B}_{v}({\bmi x};k_{x},k_{y})}\right\rangle+
\left\langle{{\bmi B}_{u}({\bmi x};-k_{x},-k_{y}),{ \partial 
}_{z}{\bmi B}_{v}({\bmi x};k_{x},k_{y})}\right\rangle \nonumber\\
&&\!\!\!\!\!\!
=
\left\langle{\rf{\bmi A}(z,x;{\partial }_{x},-{k}_{y})
{\bmi B}_{u}({\bmi x};-k_{x},-k_{y}),
{\bmi B}_{v}({\bmi x};k_{x},k_{y})}\right\rangle\nonumber\\
&&\qquad\qquad
+
\left\langle{{\bmi B}_{u}({\bmi x};-k_{x},-k_{y}),
\rf{\bmi A}(z,x;{\partial }_{x},{k}_{y}){\bmi B}_{v}
({\bmi x};k_{x},k_{y})}\right\rangle={\bf 
0}
\end{eqnarray}
because 
\begin{equation} \label{rfAsym}
\left\langle{\rf{\bmi A}(z;{\partial }_{x},-{k}_{y})\cdot,\cdot}\right\rangle=-
\left\langle{\cdot,\rf{\bmi A}(z;{\partial }_{x},{k}_{y})\cdot}\right\rangle
\end{equation}
which follows after two operations of integration by parts.



%%%%%%%%%%%%%%%%%%%%%%%%%%%%%%%%%%%%%%%%%%%%%%%%%%%%%%%%%%%%%%%%%%%%%%%%%%%%%%%%%%
A final property of the Jost solutions  ${\bmi  B}_{u}$ and ${\bmi  
B}_{v}$ is the {\em closure condition\/} ({\em vid.\/} {\sc  Box}~$\bf 7.1$)
\begin{eqnarray} \label{closure}
\lefteqn{{\bmi  T}({ x};k_y) \equiv
\frac{1}{2 \pi}\int_{-\infty}^{\infty} \left[{\bmi B}_{ v}({\bmi x};k_{x},k_{y}) \langle{  
{\bmi B}_{u}({\bmi x};-k_{x},-k_{y}),  {\bmi B}_{v}({\bmi x};k_{x},k_{y})
}\rangle^{-1} \right.}\nonumber\\
&&\qquad\qquad\qquad\qquad\qquad\qquad\times \int_{-\infty}^{\infty} \langle{  
{\bmi B}_{u}({\bmi x};-k_{x},-k_{y}),  {\bmi  T}({ x};k_y)
}\rangle {\rm d}x \nonumber\\
&&- 
{\bmi B}_{ u} ({\bmi x};k_{x},k_{y}) \langle{  
{\bmi B}_{u}({\bmi x};-k_{x},-k_{y}),  
{\bmi B}_{v}({\bmi x};k_{x},k_{y})
}\rangle^{-T}\nonumber\\
&&\qquad\qquad\qquad\qquad\qquad\qquad\times \left.
\int_{-\infty}^{\infty} \langle{ {\bmi B}_{v}({\bmi x};-k_{x},-k_{y}),  
{\bmi  T}({ x};k_y)}\rangle {\rm d}x\right]{\rm d}k_x
\end{eqnarray}
for ${\bmi  T}$ in the space of {\em $6-$component,  Fourier transformable 
vector fields\/}.
The closure condition (\ref{closure}), 
will be needed in the construction of the time harmonic elastic wave 
Green's tensor for a stratified half space.
%%%%%%%%%%%%%%%%%%%%%%%%%%%%%%%%%%%%%%%%%%%%%%%%%%%%%%%%%%%%%%%%%%%%%%%%%%%%%%%%%%%%%
%%%%%%%%%%%%%%%%%%%%%%%%%%%%%%%%%%%%%%%%%%%%%%%%%%%%%%%%%%%%%%%%%%%%%%%%%%%%%%%%%%%%%
\begin{figure}
\fbox{
\begin{minipage} {142.0mm}
{\sc \bf  Box} $\bf 7.1$ {\sc The closure property of the elastodynamic 
Jost solutions in a basis of plane waves} \newline
\vspace{4mm}

For convenience, we introduce the short hand notation
\begin{equation} \label{tildeBdef}
\left[{\bmi  B}_{u}({\bmi x}; k_x,k_y,\omega),
{\bmi  B}_{v}({\bmi x}; k_x,k_y,\omega)\right]=
{\rm e}^{i k_x x}
\left[\tilde{\bmi  B}_{u}(k_x,k_y),\tilde{\bmi  B}_{v}(k_x,k_y)\right]
\end{equation}
where $\left[{\bmi  B}_{u}({\bmi x}; k_x,k_y,\omega),
{\bmi  B}_{v}({\bmi x}; k_x,k_y,\omega)\right]$ was defined in (\ref{fundamental}).

$\quad$ Using the bilinear form (\ref{functional}), we define the {\em Gramm\/} or 
{\em overlap\/} matrix 
$\langle \tilde{\bmi  B}(-k_x,-k_y), \tilde{\bmi  B}(k_x,k_y) \rangle$, 
of  $[\tilde{\bmi  B}_{u},\tilde{\bmi  B}_{v}]$ from 
(\ref{tildeBdef}) according to
\begin{eqnarray*}
\lefteqn{\!\!\!
\left\langle{ \tilde{\bmi  B}(-k_x,-k_y),\tilde{\bmi  B}(k_x,k_y)}\right\rangle 
\!=\!\left[{\begin{array}{cc}\left\langle{\tilde{\bmi  B}_{u}(-k_x,-k_y),\tilde{\bmi 
B}_{u}(k_x,k_y)}\right\rangle&
\left\langle{\tilde{\bmi  B}_{u}(-k_x,-k_y),\tilde{\bmi  B}_{v}(k_x,k_y)}\right\rangle\\ 
\left\langle{\tilde{\bmi  B}_{v}(-k_x,-k_y),\tilde{\bmi  B}_{u}(k_x,k_y)}\right\rangle&
\left\langle{\tilde{\bmi 
B}_{v}(-k_x,-k_y),\tilde{\bmi  B}_{v}(k_x,k_y)}
\right\rangle\end{array}}\right]}\\
&&=\left({\begin{array}{cc}{\bf 
0}_{3}&\left\langle{
\tilde{\bmi  B}_{u}(-k_x,-k_y),\tilde{\bmi  B}_{v}(k_x,k_y)}\right\rangle\\ 
-{\left\langle{\tilde{\bmi  B}_{u}(-k_x,-k_y),\tilde{\bmi  B}_{v}(k_x,k_y)}\right\rangle}^{T}&{\bf 
0}_{3}\end{array}}\right)
\end{eqnarray*} 
but $\tilde{\bmi  B}(k_x,k_y){\left\langle{\tilde{\bmi  
B}(-k_x,-k_y),\tilde{\bmi 
B}(k_x,k_y)}\right\rangle}^{-1}\left\langle{\tilde{\bmi  B}(-k_x,-k_y),\cdot 
}\right\rangle$ is a  {\em projection\/} onto the space spanned by  $\tilde{\bmi 
B}(k_x,k_y)$, because 
\begin{eqnarray*}
\lefteqn{
\tilde{\bmi  B}(k_x,k_y){\left\langle{\tilde{\bmi  B}(-k_x,-k_y),\tilde{\bmi  B}(k_x,k_y)}
\right\rangle}^{-1}}\\
&&\qquad \qquad \times \left\langle{\tilde{\bmi 
B}(-k_x,-k_y),\tilde{\bmi  B}(k_x,k_y){\left\langle{\tilde{\bmi  B}(-k_x,-k_y),
\tilde{\bmi  B}(k_x,k_y)}\right\rangle}^{-1}
\left\langle{\tilde{\bmi  B}(-k_x,-k_y),\cdot 
}\right\rangle}\right\rangle=\\
&&\!\!\!\!\!\!
\tilde{\bmi  B}(k_x,k_y){\left\langle{\tilde{\bmi  B}(-k_x,-k_y),\tilde{\bmi 
B}(k_x,k_y)}\right\rangle}^{-1}\left\langle{\tilde{\bmi  B}(-k_x,-k_y),
\tilde{\bmi  B}(k_x,k_y)}\right
\rangle\\
&&\qquad \qquad \times {\left\langle{\tilde{\bmi  B}(-k_x,-k_y),\tilde{\bmi  B}(k_x,k_y)}\right\rangle}^{-1}
\left\langle{\tilde{\bmi 
B}(-k_x,-k_y),\cdot 
}\right\rangle=\\
&&\!\!\!\!\!\!
\tilde{\bmi  B}(k_x,k_y){\left\langle{\tilde{\bmi  B}(-k_x,-k_y),\tilde{\bmi 
B}(k_x,k_y)}\right\rangle}^{-1}\left\langle{\tilde{\bmi  B}(k_x,k_y),\cdot 
}\right\rangle
\end{eqnarray*} 
so that 
\[
\tilde{\bmi  B}(k_x,k_y){\left\langle{\tilde{\bmi  B}(-k_x,-k_y),\tilde{\bmi 
B}(k_x,k_y)}\right\rangle}^{-1}\left\langle{\tilde{\bmi  B}(-k_x,-k_y),\cdot }
\right\rangle 
\] 
is {\em idempotent\/}.

Now, assuming we are given a field ${\bmi  T}(k_x,k_y)=\tilde{\bmi  B}(k_x,k_y)\cdot {\bmi 
C}$, with arbitrary expansion coefficients $\bmi  C$, we have
\begin{eqnarray*}
\lefteqn{ 
\tilde{\bmi  B}(k_x,k_y){\left\langle{\tilde{\bmi  B}(-k_x,-k_y),\tilde{\bmi  B}(k_x,k_y)}
\right\rangle}^{-1}\left\langle{\tilde{\bmi 
B}(-k_x,-k_y),{\bmi  T}(k_x,k_y)}\right\rangle=}\\
&&
\tilde{\bmi  B}(k_x,k_y){\left\langle{\tilde{\bmi  B}(-k_x,-k_y),\tilde{\bmi  B}(k_x,k_y)}
\right\rangle}^{-1}\left\langle{\tilde{\bmi 
B}(-k_x,-k_y),\tilde{\bmi  B}(k_x,k_y)\cdot 
 {\bmi  C}}\right\rangle=\\
&&\tilde{\bmi  B}(k_x,k_y){\left\langle{\tilde{\bmi  B}(-k_x,-k_y),
 \tilde{\bmi  B}(k_x,k_y)}\right\rangle}^{-1}
 \left\langle{\tilde{\bmi  B}(-k_x,-k_y),\tilde{\bmi  B}(k_x,k_y)}\right\rangle\cdot  
 {\bmi  C}=\\
&&\tilde{\bmi  B}(k_x,k_y)\cdot  
 {\bmi  C}={\bmi  T}(k_x,k_y)
\end{eqnarray*} 
\begin{eqnarray*}
\lefteqn{ 
{\left\langle{{\bmi  B }(-k_x,-k_y),{ \bmi  B }(k_x,k_y)}\right\rangle}^{\rm -1} 
}\\
&&={\left({\begin{array}{cc}{\bf 0}_{3}&\left\langle{\tilde{\bmi  B}_{u}(-k_x,-k_y),
\tilde{\bmi  B}_{v}(k_x,k_y)}\right\rangle\\ 
-{\left\langle{\tilde{\bmi  B}_{u}(-k_x,-k_y),\tilde{\bmi  B}_{v}
(k_x,k_y)}\right\rangle}^{T}&{\bf 0}_{3}\end{array}}\right)}^{-1}\\
&&=
\left({\begin{array}{cc}{\bf 0}_{3}&-{\left\langle{
\tilde{\bmi  B}_{u}(-k_x,-k_y),\tilde{\bmi  B}_{v}(k_x,k_y)}\right\rangle}^{-T}\\ 
{\left\langle{\tilde{\bmi  B}_{u}(-k_x,-k_y),\tilde{\bmi  B}_{v}(k_x,k_y)}\right\rangle}^{-1}
&{\bf 0}_{3}\end{array}}\right)
\end{eqnarray*} 

\vspace{6mm}
\end{minipage}
}
\end{figure}
\begin{figure}
\fbox{
\begin{minipage} {142.0mm}
{\sc \bf  Box} $\bf 7.1$ {\em continued}\newline
and using 
in the previous equation, we get
\begin{eqnarray} \label{reconstruct0}
\lefteqn{\tilde{\bmi  B}_{ v}(k_x,k_y) \langle{ 
\tilde{\bmi  B}_{u}(-k_x,-k_y),  \tilde{\bmi  B}_{v}(k_x,k_y)
}\rangle^{-1}  \langle{ 
\tilde{\bmi  B}_{u}(-k_x,-k_y),  {\bmi  T}(k_x,k_y)
}\rangle  - }\nonumber\\
&&
\tilde{\bmi  B}_{ u}(k_x,k_y) \langle{ 
\tilde{\bmi  B}_{u}(-k_x,-k_y),  \tilde{\bmi  B}_{v}(k_x,k_y)
}\rangle^{-T}
\langle{  
\tilde{\bmi  B}_{v}(-k_x,-k_y),  {\bmi  T}(k_x,k_y)
}\rangle\equiv {\bmi  T}(k_x,k_y)
\end{eqnarray}

But, using a Fourier transform of $x$ we have 
\[
{\bmi  T}(k_x,k_y)=\int_{-\infty}^{\infty} {\rm d}x\, {\rm e}^{-i k_x x} {\bmi  T}(x,k_y)
\]
and substituting this result in (\ref{reconstruct0}), gives
\begin{eqnarray} \label{reconstruct1}
\lefteqn{\tilde{\bmi  B}_{ v}(k_x,k_y) \langle{ 
\tilde{\bmi  B}_{u}(-k_x,-k_y),  \tilde{\bmi  B}_{v}(k_x,k_y)
}\rangle^{-1}  \int_{-\infty}^{\infty} \langle{ 
{\rm e}^{-i k_x x} \tilde{\bmi  B}_{u}(-k_x,-k_y),  {\bmi  T}(x,k_y)
}\rangle {\rm d}x - }\nonumber\\
&&\!\!\!\!\!\!
\tilde{\bmi  B}_{ u}(k_x,k_y) \langle{ 
\tilde{\bmi  B}_{u}(-k_x,-k_y),  \tilde{\bmi  B}_{v}(k_x,k_y)
}\rangle^{-T}
\int_{-\infty}^{\infty} \langle{ 
{\rm e}^{-i k_x x} \tilde{\bmi  B}_{v}(-k_x,-k_y),  {\bmi  T}(x,k_y)
}\rangle {\rm d}x \nonumber\\
&&\!\!\!\!\!\!\equiv {\bmi  T}(k_x,k_y)
\end{eqnarray}
However,
\[
{\bmi  T}(x,k_y)=\frac{1}{2 \pi}\int_{-\infty}^{\infty} {\rm d}k\, 
{\rm e}^{i k_x x} {\bmi  T}(k_x,k_y)
\]
so that operating on both sides of  (\ref{reconstruct1}) with 
$\frac{1}{2 \pi}\int_{-\infty}^{\infty} {\rm d}k\, 
{\rm e}^{i k_x x}$ gives
\begin{eqnarray} \label{reconstruct2}
\lefteqn{
\!\!\!\!\!\!
\frac{1}{2 \pi}\int_{-\infty}^{\infty}\!\!\!  \left[{\rm e}^{i k_x x}\tilde{\bmi  B}_{ v}(k_x,k_y) \langle{ 
\tilde{\bmi  B}_{u}(-k_x,-k_y),  \tilde{\bmi  B}_{v}(k_x,k_y)
}\rangle^{-1}\!\!\!  \int_{-\infty}^{\infty}\!\!\! \langle{ 
{\rm e}^{-i k_x x} \tilde{\bmi  B}_{u}(-k_x,-k_y),  {\bmi  T}(x,k_y)
}\rangle\, {\rm d}x  \right.}\nonumber\\
&&\!\!\!\!\!\!
\left. -{\rm e}^{i k_x x} \tilde{\bmi  B}_{ u}(k_x,k_y) \langle{ 
\tilde{\bmi  B}_{u}(-k_x,-k_y),  \tilde{\bmi  B}_{v}(k_x,k_y)
}\rangle^{-T}\!\!\!
\int_{-\infty}^{\infty} \!\!\! \langle{ 
{\rm e}^{-i k_x x} \tilde{\bmi  B}_{v}(-k_x,-k_y),  {\bmi  T}(x,k_y)
}\rangle {\rm d}x
\right]{\rm d}k\nonumber\\
&&\!\!\!\!\!\!\equiv {\bmi  T}(x,k_y)
\end{eqnarray}

Using the definition (\ref{tildeBdef}) in (\ref{reconstruct2}) we get
\begin{eqnarray} \label{reconstruct3}
\lefteqn{
\frac{1}{2 \pi}\int_{-\infty}^{\infty} {\rm d}k\, \bigg[
{\bmi  B}_{ v}({\bmi x};k_x,k_y,\omega) \langle{ 
{\bmi  B}_{u}({\bmi x};-k_x,-k_y,\omega),  {\bmi  B}_{v}({\bmi x};k_x,k_y,\omega)
}\rangle^{-1}    }\nonumber\\
&&\qquad\qquad\qquad\qquad\qquad\qquad\qquad\qquad\times \int_{-\infty}^{\infty} \langle{ 
{\bmi  B}_{u}({\bmi x};-k_x,-k_y,\omega),  {\bmi  T}(x,k_y)
}\rangle {\rm d}x\nonumber\\
&&
-{\bmi  B}_{ u}({\bmi x};k_x,k_y,\omega) \langle{ 
{\bmi  B}_{u}({\bmi x};-k_x,-k_y,\omega),  {\bmi  B}_{v}({\bmi x};k_x,k_y,\omega)
}\rangle^{-T}\nonumber\\
&& \qquad\qquad\qquad\qquad\qquad\qquad\qquad\qquad\times
 \int_{-\infty}^{\infty} \langle{ 
{\bmi  B}_{v}({\bmi x};-k_x,-k_y,\omega),  {\bmi  T}(x,k_y)
}\rangle {\rm d}x
\bigg]\nonumber\\
&&\equiv {\bmi  T}(x,k_y)
\end{eqnarray}

\vspace{6mm}
\end{minipage}
}
\end{figure}

%%%%%%%%%%%%%%%%%%%%%%%%%%%%%%%%%%%%%%%%%%%%%%%%%%%%%%%%%%%%%%%%%%%%%%%%%%%%%%%%%%%%%
%%%%%%%%%%%%%%%%%%%%%%%%%%%%%%%%%%%%%%%%%%%%%%%%%%%%%%%%%%%%%%%%%%%%%%%%%%%%%%%%%%%%%


\newpage
\section{A Jost solution construction for the Lippmann-Schwinger equation} 
\label{sectionstratG}
In this section, we  re-derive the {\em  Lippmann-Schwinger equation\/} of 
(\ref{LSmat}), but this time we employ the {\em Jost solutions\/} of (\ref{fundamental}).

The {\em  heterogeneous half space Green's tensor field \/}  
${\bmi G}({\bmi  x} ; {\bmi  x}_{ S};k_y)$, 
driven by a point source at ${\bmi  
x}_S\equiv (z_S,x_S)$,  obeys the
following  boundary value problem, defined in terms of the 
bilinear form $\langle \cdot, \cdot \rangle$ of (\ref{functional}) as:
\begin{samepage} 
\begin{eqnarray} \label{Geqntotal} 
\!\!\!\!\!\!
\lefteqn{
\left\{ \begin{array}{l}
\int_{-\infty}^{\infty} \langle {\bmi B}_{ v} ((f+,x);-k_x,-k_y)
, {\bmi G}  ({\bmi  x} ; {\bmi  x}_{ S}; k_y)\rangle \, {\rm d}x = 
{\bf 0}_{ 3   \times  6} 
\\
 \left[{\partial}_{z} -  
\rf{\bmi A}(z,{\partial}_{x}, k_y)\right] {\bmi G}   ({\bmi  x} ; {\bmi  
x}_{ S}; k_y)= 
{\mit \Delta}{\bmi A}({\bmi x},{\partial}_{x}, k_y) 
{\bmi G}   ({\bmi  x} ; {\bmi  x}_{ S}; k_y) +  
 \delta({\bmi  x} - {\bmi  x}_{ S}) {\bmi  I}_{6} \\
\int_{-\infty}^{\infty} \langle {\bmi B}_{ u} ((L+,x);-k_x,-k_y), 
{\bmi G}  ({\bmi  x} ; {\bmi  x}_{ S}; k_y)\rangle \, {\rm d}x
= {\bf 0}_{3   \times  6} 
\end{array} \right.  
}\nonumber\\ 
&&
\end{eqnarray}
\end{samepage}
where $\rf{\bmi A}(z,{\partial}_{x}, k_y)$ was given by (\ref{Adefref}) and 
${\mit \Delta}{\bmi A}({\bmi x},{\partial}_{x}, k_y)$ by (\ref{dAdef}).
 The $3 $ boundary conditions 
\[
\int_{-\infty}^{\infty} \langle {\bmi B}_{ v} ((f+,x);-k_x,-k_y)
, {\bmi G}  ({\bmi  x} ; {\bmi  x}_{ S};k_y)\rangle \, {\rm d}x = 
{\bf 0}_{ 3   \times  6}
\]
 ensure that the traction field vanishes at the 
free surface, while the $3 $ boundary conditions 
\[
\int_{-\infty}^{\infty} \langle {\bmi B}_{ u} ((L+,x);-k_x,-k_y), 
{\bmi G}  ({\bmi  x} ; {\bmi  x}_{ S};k_y)\rangle \, {\rm d}x
= {\bf 0}_{3   \times  6}
\]
ensure that
there is no $\uparrow\! \!-going$ incident field at the base of the 
stratification. 

%%%%%%%%%%%%%%%%%%%%%%%%%%%%%
We now construct the depth evolution equations for the bilinear functional 
\[
\int_{-\infty}^{\infty} \langle{\bmi B}_{u}({\bmi x};-k_x,-k_y), 
{\bmi G}({\bmi  x} ; {\bmi  x}_{ S};k_y)\ \rangle \, {\rm d}x
\]
 comprised
of ${\bmi B}_{u}({\bmi x};-k_x,-k_y)$ from (\ref{Budef})  
and the  Greens tensor ${\bmi G}({\bmi  x} ; {\bmi  x}_{ S};k_y)$ of 
(\ref{Geqntotal}) {\em viz. \/},
\begin{eqnarray} \label{(2.4.2)} 
\lefteqn{{\partial }_{z} \int_{-\infty}^{\infty} \left\langle{  {\bmi B}_{u}({\bmi x};-k_x,-k_y),  
{\bmi G}({\bmi  x} ; {\bmi  x}_{ S};k_y)
}\right\rangle \,{\rm d}x  =
}\nonumber\\
&&
\int_{-\infty}^{\infty} \left\langle{ {\partial }_{z}  {\bmi B}_{u}({\bmi x};-k_x,-k_y),  
{\bmi G}({\bmi  x} ; {\bmi  x}_{ S};k_y)
}\right\rangle \,{\rm d}x \nonumber\\
&&+
\int_{-\infty}^{\infty} \left\langle{   {\bmi B}_{u}({\bmi x};-k_x,-k_y),  
{\partial }_{z} {\bmi G}({\bmi  x} ; {\bmi  x}_{ S};k_y)
}\right\rangle \,{\rm d}x=\nonumber\\
&&
\int_{-\infty}^{\infty} \bigg[ \langle \rf{\bmi A}({z},{\partial}_{x}, -k_y) {\bmi B}_{u}({\bmi x};-k_x,-k_y),  
{\bmi G}({\bmi  x} ; {\bmi  x}_{ S};k_y) \rangle   \nonumber\\
&&
+\langle{\bmi B}_{u}({\bmi x};-k_x,-k_y), 
\rf{\bmi A}({z},{\partial}_{x}, k_y) {\bmi G}({\bmi  x} ; {\bmi  x}_{ S};k_y) 
\rangle\bigg]\,{\rm d}x\nonumber\\
&&
+ \delta (z-{z}_{S}) {\bmi B}{}_{u}^{T}({\bmi x}_{S};-k_x,-k_y)
{\bmi  J}_{6}  \nonumber\\
&&
+ 
\int_{-\infty}^{\infty} \langle{\bmi B}_{u}({\bmi x};-k_x,-k_y), {\mit \Delta}
{\bmi A}({\bmi x},{\partial}_{x}, k_y) {\bmi G}({\bmi  x} ; {\bmi  x}_{ S};k_y) \rangle \,{\rm d}x
\end{eqnarray} 
but
\begin{eqnarray} \label{Asym1}
\lefteqn{\int_{-\infty}^{\infty} \bigg[ \langle \rf{\bmi A}({z},{\partial}_{x}, -k_y) {\bmi B}_{u}({\bmi x};-k_x,-k_y),  
{\bmi G}({\bmi  x} ; {\bmi  x}_{ S};k_y) \rangle }  \nonumber\\
&&
\qquad\qquad\quad
+\langle{\bmi B}_{u}({\bmi x};-k_x,-k_y), 
\rf{\bmi A}({z},{\partial}_{x}, k_y) {\bmi G}({\bmi  x} ; {\bmi  x}_{ S};k_y) 
\rangle\bigg]\,{\rm d}x= {\bf 0}_3
\end{eqnarray}
from (\ref{rfAsym}),
due to the {\em quasi-Hamiltonian symmetry\/} of $\rf{\bmi A}$. Using 
(\ref{Asym1}) in (\ref{(2.4.2)}) gives
\begin{eqnarray} \label{BuG} 
\lefteqn{{\partial }_{z} \int_{-\infty}^{\infty} \left\langle{  {\bmi B}_{u}({\bmi x};-k_x,-k_y),  
{\bmi G}({\bmi  x} ; {\bmi  x}_{ S};k_y)
}\right\rangle \,{\rm d}x  =\delta (z-{z}_{S}) 
{\bmi B}{}_{u}^{T}({\bmi x}_{S};-k_x,-k_y)
{\bmi  J}_{6}
}\nonumber\\
&&
\qquad\qquad\qquad+ 
\int_{-\infty}^{\infty} \langle{\bmi B}_{u}({\bmi x};-k_x,-k_y), {\mit \Delta}
{\bmi A}({\bmi x},{\partial}_{x}, k_y) {\bmi G}({\bmi  x} ; {\bmi  x}_{ S};k_y) \rangle \,{\rm d}x
\end{eqnarray} 




Upon integrating  (\ref{BuG}) over the depth interval 
$z+ \le z \le {z}_{L}+$ we get 
\begin{eqnarray*}  
\lefteqn{\overbrace{\left.\int_{-\infty}^{\infty} \left\langle{  {\bmi B}_{u}({\bmi x};-k_x,-k_y),  
{\bmi G}({\bmi  x} ; {\bmi  x}_{ S};k_y)
}\right\rangle \,{\rm d}x \right|_{{z}_{L+}}}^{={\bf 0} \ \mbox{by 
radiation conditions}}}\\
&&\qquad\qquad\qquad
-\left.\int_{-\infty}^{\infty} \left\langle{  {\bmi B}_{u}({\bmi x};-k_x,-k_y),  
{\bmi G}({\bmi  x} ; {\bmi  x}_{ S};k_y)
}\right\rangle \,{\rm d}x \right|_{{z}_{+}}=\\
&&
H({z}_{S}-z) {\bmi B}{}_{u}^{T}({\bmi x}_{S};-k_x,-k_y)
{\bmi  J}_{6}  \\
&&\qquad\qquad\qquad
+\int_{z}^{z_L}
\int_{-\infty}^{\infty} \langle{\bmi B}_{u}({\bmi x};-k_x,-k_y), {\mit \Delta}
{\bmi A}({\bmi x},{\partial}_{x}, k_y) {\bmi G}({\bmi  x} ; {\bmi  x}_{ S};k_y) \rangle 
\,{\rm d}x\,{\rm d}z
\end{eqnarray*} 
so that
\begin{eqnarray}  \label{(2.4.3)}
\lefteqn{\!\!\!\!\!\!
\left.\int_{-\infty}^{\infty} \left\langle{  {\bmi B}_{u}({\bmi x};-k_x,-k_y),  
{\bmi G}({\bmi  x} ; {\bmi  x}_{ S};k_y)
}\right\rangle \,{\rm d}x \right|_{{z}_{+}}=-H({z}_{S}-z) {\bmi B}{}_{u}^{T}({\bmi x}_{S};-k_x,-k_y)
{\bmi  J}_{6}}\nonumber\\
&&\qquad
-\int_{z}^{z_L}
\int_{-\infty}^{\infty} \langle{\bmi B}_{u}({\bmi x};-k_x,-k_y), {\mit \Delta}
{\bmi A}({\bmi x},{\partial}_{x}, k_y) {\bmi G}({\bmi  x} ; {\bmi  x}_{ S};k_y) \rangle 
\,{\rm d}x\,{\rm d}z
\end{eqnarray} 


%%%%%%%%%%%%%%%%%%%%%%%%%%%%%%%%%%%%%%%%%%%%%%%%%%%%%%%%%%%%%%%%%%%%%%%%%%%
%%%%%%%%%%%%%%%%%%%%%%%%%%%%%%%%%%%%%%%%%%%%%%%%%%%%%%%%%%%%%%%%%%%%%%%%%%%
%%%%%%%%%%%%%%%%%%%%%%%%%%%%%%%%%%%%%%%%%%%%%%%%%%%%%%%%%%%%%%%%%%%%%%%%%%%
%%%%%%%%%%%%%%%%%%%%%%%%%%%%%%%%%%%%%%%%%%%%%%%%%%%%%%%%%%%%%%%%%%%%%%%%%%%
%%%%%%%%%%%%%%%%%%%%%%%%%%%%%%%%%%%%%%%%%%%%%%%%%%%%%%%%%%%%%%%%%%%%%%%%%%%
Similarly, we construct the evolution equations for 
\[
\int_{-\infty}^{\infty} \langle{\bmi B}_{v}({\bmi x};-k_x,-k_y), 
{\bmi G}({\bmi  x} ; {\bmi  x}_{ S};k_y)\ \rangle \, {\rm d}x
\]
 comprised
of ${\bmi B}_{v}({\bmi x};-k_x,-k_y)$ from (\ref{Bvdef})  
and the  Greens tensor ${\bmi G}({\bmi  x} ; {\bmi  x}_{ S};k_y)$ of 
(\ref{Geqntotal}) {\em viz. \/},
\begin{eqnarray} \label{(2.4.4)} 
\lefteqn{{\partial }_{z} \int_{-\infty}^{\infty} \left\langle{  {\bmi B}_{v}({\bmi x};-k_x,-k_y),  
{\bmi G}({\bmi  x} ; {\bmi  x}_{ S};k_y)
}\right\rangle \,{\rm d}x  =
}\nonumber\\
&&
\int_{-\infty}^{\infty} \left\langle{ {\partial }_{z}  {\bmi B}_{v}({\bmi x};-k_x,-k_y),  
{\bmi G}({\bmi  x} ; {\bmi  x}_{ S};k_y)
}\right\rangle \,{\rm d}x \nonumber\\
&&+
\int_{-\infty}^{\infty} \left\langle{   {\bmi B}_{v}({\bmi x};-k_x,-k_y),  
{\partial }_{z} {\bmi G}({\bmi  x} ; {\bmi  x}_{ S};k_y)
}\right\rangle \,{\rm d}x=\nonumber\\
&&
\int_{-\infty}^{\infty} \bigg[ \langle \rf{\bmi A}({z},{\partial}_{x}, -k_y) {\bmi B}_{v}({\bmi x};-k_x,-k_y),  
{\bmi G}({\bmi  x} ; {\bmi  x}_{ S};k_y) \rangle   \nonumber\\
&&
+\langle{\bmi B}_{v}({\bmi x};-k_x,-k_y), 
\rf{\bmi A}({z},{\partial}_{x}, k_y) {\bmi G}({\bmi  x} ; {\bmi  x}_{ S};k_y) 
\rangle\bigg]\,{\rm d}x\nonumber\\
&&
+ \delta (z-{z}_{S}) {\bmi B}{}_{v}^{T}({\bmi x}_{S};-k_x,-k_y)
{\bmi  J}_{6}  \nonumber\\
&&
+ 
\int_{-\infty}^{\infty} \langle{\bmi B}_{v}({\bmi x};-k_x,-k_y), {\mit \Delta}
{\bmi A}({\bmi x},{\partial}_{x}, k_y) {\bmi G}({\bmi  x} ; {\bmi  x}_{ S};k_y) \rangle \,{\rm d}x
\end{eqnarray} 
but
\begin{eqnarray} \label{Asym2}
\lefteqn{\int_{-\infty}^{\infty} \bigg[ \langle \rf{\bmi A}({z},{\partial}_{x}, -k_y) {\bmi B}_{v}({\bmi x};-k_x,-k_y),  
{\bmi G}({\bmi  x} ; {\bmi  x}_{ S};k_y) \rangle }  \nonumber\\
&&
\qquad\qquad\quad
+\langle{\bmi B}_{v}({\bmi x};-k_x,-k_y), 
\rf{\bmi A}({z},{\partial}_{x}, k_y) {\bmi G}({\bmi  x} ; {\bmi  x}_{ S};k_y) 
\rangle\bigg]\,{\rm d}x= {\bf 0}_3
\end{eqnarray}
from (\ref{rfAsym}),
due to the {\em quasi-Hamiltonian symmetry\/} of $\rf{\bmi A}$. Using 
(\ref{Asym2}) in (\ref{(2.4.4)}) gives
\begin{eqnarray} \label{BvG} 
\lefteqn{{\partial }_{z} \int_{-\infty}^{\infty} \left\langle{  {\bmi B}_{v}({\bmi x};-k_x,-k_y),  
{\bmi G}({\bmi  x} ; {\bmi  x}_{ S};k_y)
}\right\rangle \,{\rm d}x  =\delta (z-{z}_{S}) 
{\bmi B}{}_{v}^{T}({\bmi x}_{S};-k_x,-k_y)
{\bmi  J}_{6}
}\nonumber\\
&&
\qquad\qquad\qquad+ 
\int_{-\infty}^{\infty} \langle{\bmi B}_{v}({\bmi x};-k_x,-k_y), {\mit \Delta}
{\bmi A}({\bmi x},{\partial}_{x}, k_y) {\bmi G}({\bmi  x} ; {\bmi  x}_{ S};k_y) \rangle \,{\rm d}x
\end{eqnarray} 




Upon integrating  (\ref{BvG}) over the depth interval 
$f+ \le z \le z-$ we get 
\begin{eqnarray*}  
\lefteqn{{\left.\int_{-\infty}^{\infty} \left\langle{  {\bmi B}_{v}({\bmi x};-k_x,-k_y),  
{\bmi G}({\bmi  x} ; {\bmi  x}_{ S};k_y)
}\right\rangle \,{\rm d}x \right|_{{z}-}}}\\
&&\qquad\qquad\qquad
-\overbrace{\left.\int_{-\infty}^{\infty} \left\langle{  {\bmi B}_{v}({\bmi x};-k_x,-k_y),  
{\bmi G}({\bmi  x} ; {\bmi  x}_{ S};k_y)
}\right\rangle \,{\rm d}x \right|_{{z}_{f}+}}^{={\bf 0} \ \mbox{by 
free surface conditions}}=\\
&&
H(z-{z}_{S}) {\bmi B}{}_{v}^{T}({\bmi x}_{S};-k_x,-k_y)
{\bmi  J}_{6}  \\
&&\qquad\qquad\qquad
+\int_{z_f}^{z}
\int_{-\infty}^{\infty} \langle{\bmi B}_{v}({\bmi x};-k_x,-k_y), {\mit \Delta}
{\bmi A}({\bmi x},{\partial}_{x}, k_y) {\bmi G}({\bmi  x} ; {\bmi  x}_{ S};k_y) \rangle 
\,{\rm d}x\,{\rm d}z
\end{eqnarray*} 
so that
\begin{eqnarray}  \label{(2.4.5)}
\lefteqn{\!\!\!\!\!\!
\left.\int_{-\infty}^{\infty} \left\langle{  {\bmi B}_{v}({\bmi x};-k_x,-k_y),  
{\bmi G}({\bmi  x} ; {\bmi  x}_{ S};k_y)
}\right\rangle \,{\rm d}x \right|_{{z}_{-}}=H(z-{z}_{S}) {\bmi B}{}_{v}^{T}({\bmi x}_{S};-k_x,-k_y)
{\bmi  J}_{6}}\nonumber\\
&&\qquad
+\int_{z_f}^{z}
\int_{-\infty}^{\infty} \langle{\bmi B}_{v}({\bmi x};-k_x,-k_y), {\mit \Delta}
{\bmi A}({\bmi x},{\partial}_{x}, k_y) {\bmi G}({\bmi  x} ; {\bmi  x}_{ S};k_y) \rangle 
\,{\rm d}x\,{\rm d}z
\end{eqnarray} 
%%%%%%%%%%%%%%%%%%%%%%%%%%%%%%%%%%%%%%%%%%%%%%%%%%%%%%%%%%%%%%%%%%%%%%%%%
%%%%%%%%%%%%%%%%%%%%%%%%%%%%%%%%%%%%%%%%%%%%%%%%%%%%%%%%%%%%%%%%%%%%%%%%%
%%%%%%%%%%%%%%%%%%%%%%%%%%%%%%%%%%%%%%%%%%%%%%%%%%%%%%%%%%%%%%%%%%%%%%%%%
Next we form the  linear combination 
\begin{eqnarray} \label{(2.4.6)}
\lefteqn{\frac{1}{2 \pi}\int_{-\infty}^{\infty} \bigg[{\bmi B}_{ v}({\bmi x};k_{x},k_{y}) \langle{  
{\bmi B}_{u}({\bmi x};-k_{x},-k_{y}),  {\bmi B}_{v}({\bmi x};k_{x},k_{y})
}\rangle^{-1}\times (\ref{(2.4.3)}) }\nonumber\\
&&- 
{\bmi B}_{ u} ({\bmi x};k_{x},k_{y}) \langle{  
{\bmi B}_{u}({\bmi x};-k_{x},-k_{y}),  
{\bmi B}_{v}({\bmi x};k_{x},k_{y})
}\rangle^{-T}\times (\ref{(2.4.5)})\bigg]\,{\rm d}k_x
\end{eqnarray}
concocted so that the resulting left hand side reduces to

\begin{samepage}
\begin{eqnarray*}
\lefteqn{{\bmi  G}({\bmi x},{\bmi x}_{S};k_y) \equiv
\frac{1}{2 \pi}\int_{-\infty}^{\infty} \left[
{\bmi B}_{ v}({\bmi x};k_{x},k_{y}) \langle{  
{\bmi B}_{u}({\bmi x};-k_{x},-k_{y}),  {\bmi B}_{v}({\bmi x};k_{x},k_{y})
}\rangle^{-1} \right.}\nonumber\\
&&\qquad\qquad\qquad\qquad\qquad\times \int_{-\infty}^{\infty} \langle{  
{\bmi B}_{u}({\bmi x};-k_{x},-k_{y}),  {\bmi  G}({\bmi x},{\bmi 
x}_{S};k_y)
}\rangle {\rm d}x \nonumber\\
&&- 
{\bmi B}_{ u} ({\bmi x};k_{x},k_{y}) \langle{  
{\bmi B}_{u}({\bmi x};-k_{x},-k_{y}),  
{\bmi B}_{v}({\bmi x};k_{x},k_{y})
}\rangle^{-T}\\
&&\qquad\qquad\qquad\qquad\qquad\times \left.
\int_{-\infty}^{\infty} \langle{ {\bmi B}_{v}({\bmi x};-k_{x},-k_{y}),  
{\bmi  G}({\bmi x},{\bmi x}_{S};k_y)}\rangle {\rm d}x\right]{\rm d}k_x\nonumber
\end{eqnarray*}
\end{samepage}
which follows from the {\em closure\/} property (\ref{closure}), of the 
Jost solutions $\{ {\bmi B}_{ u},{\bmi B}_{ v} \}$. From
(\ref{(2.4.6)}) we get the {\em Lippmann-Schwinger\/} integral equation
\begin{eqnarray} \label{LSformal}
{\bmi G}({\bmi x},{\bmi x}_{S};k_y)&=&
\rf{\bmi G}\!({\bmi x},{\bmi x}_{S};k_y)\nonumber\\
&+&\int_{f+}^{L+}\!\!\!\int_{-\infty}^{\infty} 
\!\! \rf{\bmi G}\!({\bmi x},{\bmi x}^{\prime};k_y) 
{\mit \Delta}{\bmi A}({\bmi x}^{\prime},{\partial}_{x};k_y) 
{\bmi G}({\bmi x}^{\prime},{\bmi x}_{S};k_y)\, {\rm d}{\bmi x}^{\prime}
\end{eqnarray}
where the {\em stratified reference Green's tensor\/} $\rf{\bmi G} 
({\bmi x},{\bmi x}_{S};k_y)$ in (\ref{LSformal}) is defined as
\begin{eqnarray}\label{GT0}
\lefteqn{\rf{\bmi G}({\bmi x},{\bmi x}_{S};k_y)= 
\frac{1}{2 \pi}\int_{-\infty}^{\infty} {\rm d}k_x}\\
&&
\times\left\{
{\begin{array}{ll}
\textstyle
-{\bmi B}_{ v}({\bmi x};{\bmi K}_{T}) \langle{  
{\bmi B}_{u}({\bmi x}^{\prime};-{\bmi K}_{T}),  
{\bmi B}_{v}({\bmi x}^{\prime};{\bmi K}_{T})
}\rangle^{-1}
{\bmi B}^{T}_{u}({\bmi x}_{S};-{\bmi K}_{T}) {\bmi  J}_{6}& 
z<z_{S} \\
\textstyle
-{\bmi B}_{ u} ({\bmi x};{\bmi K}_{T}) \langle{  
{\bmi B}_{u}({\bmi x}^{\prime};-{\bmi K}_{T}),  
{\bmi B}_{v}({\bmi x}^{\prime};{\bmi K}_{T})
}\rangle^{-T}
{\bmi B}^{T}_{v}({\bmi x}_{S};-{\bmi K}_{T}) {\bmi  J}_{6}& z>z_{S} 
\end{array} }\right.\nonumber
\end{eqnarray} 
for an arbitrary depth point ${\bmi x}^{\prime}$, a source at
${\bmi x}_S$ a receiver at ${\bmi x}$, and we have introduced  
${\bmi K}_{T}=(k_x , k_y)$ for the horizontal component of the wavenumber.


The {\em inverse invariant\/} 
$\langle{ \ {\bmi B}_{u},  {\bmi B}_{v}\ }\rangle^{-1}$  
in  (\ref{GT0}) is singular, i.e.
\begin{eqnarray} \label{polecondition}
\lefteqn{\mbox{Det}\langle {\bmi  B}_{u}({\bmi x}^{\prime},-{\bmi K}_{T}),
{\bmi B}_{v}({\bmi x}^{\prime},{\bmi K}_{T}) \rangle
=}\nonumber\\
&&\!\!\!\!\!\!\!\!\!
\left|{\begin{array}{ccc}
\left\langle{(u,P,-{\bmi K}_{T}),(v,P,{\bmi K}_{T})}\right\rangle&
\left\langle{(u,P,-{\bmi K}_{T}),(v,S,{\bmi K}_{T})}\right\rangle&
\left\langle{(u,P,-{\bmi K}_{T}),(v,H,{\bmi K}_{T})}\right\rangle\\ 
\left\langle{(u,S,-{\bmi K}_{T}),(v,P,{\bmi K}_{T})}\right\rangle&
\left\langle{(u,S,-{\bmi K}_{T}),(v,S,{\bmi K}_{T})}\right\rangle&
\left\langle{(u,S,-{\bmi K}_{T}),(v,H,{\bmi K}_{T})}\right\rangle\\
\left\langle{(u,H,-{\bmi K}_{T}),(v,P,{\bmi K}_{T})}\right\rangle&
\left\langle{(u,H,-{\bmi K}_{T}),(v,S,{\bmi K}_{T})}\right\rangle&
\left\langle{(u,H,-{\bmi K}_{T}),(v,H,{\bmi K}_{T})}\right\rangle
\end{array}}\right|\nonumber\\
&&\!\!\!\!\!\!
=0
\end{eqnarray}
at  the surface wave poles ${\bmi K}_{T}={\bmi k}_{\mu}$ on the real axis. In the basic method, we avoid 
 the surface wave poles by moving them off the real wavenumber axis 
 through the agency of the complex frequency. 
However,  surface waves are often an important phase on seismic 
 records at regional offsets and beyond, and it would be very worthwhile 
 to explicitly extract these phases from the companion body wave phases 
 by {\em evaluating the residues\/} at the singular points of 
 $\langle{ \ {\bmi B}_{u},  {\bmi B}_{v}\ }\rangle^{-1}$.  


\newpage
%%%%%%%%%%%%%%%%%%%%%%%%%%%%%%%%%%%%%%%%%%%%%%%%%%%%%%%%%%%%%%%%%%%%%%%%%%%%%%
%%%%%%%%%%%%%%%%%%%%%%%%%%%%%%%%%%%%%%%%%%%%%%%%%%%%%%%%%%%%%%%%%%%%%%%%%%%%%%%%%%%%%
\section {Extracting  the surface waves}
%%%%%%%%%%%%%%%%%%%%%%%%%%%%%%%%%%%%%%%%%%%%%%%%%%%%%%%%%%%%%%%%%%%%%%%%%%%%%%
%%%%%%%%%%%%%%%%%%%%%%%%%%%%%%%%%%%%%%%%%%%%%%%%%%%%%%%%%%%%%%%%%%%%%%%%%%%%%%%%%%%%%



We deform  both ends of the $k_x$ wavenumber contour 
({\em vid.\/} Fig.\ref{contour}) into the $\Im(k_x) \gtlt 0$ plane for 
$(x-x_S) \gtlt 0$, so that the wave number integrand in (\ref{GT0}) 
decays exponentially.
Following {\em Kennett\/}~(1981), we account for the surface waves by 
calculating the {\em residues\/} of the integrand in (\ref{GT0})  at the 
values of ${\bmi K}_{T}={\bmi k}_{\mu}$ where the pole condition 
(\ref{polecondition}) is fulfilled.

\begin{figure}
\HideDisplacementBoxes
%\ShowDisplacementBoxes
\centerline{\BoxedEPSF{contour.EPSF scaled 800}}
\caption{The wavenumber contour to separate body wave and surface wave 
phases. Body wave phases are synthesized from the wavenumber integrals, 
and the surface wave phases from the residues at the poles, identified
as black dots. }
\protect\label{contour}
\end{figure}
 
We recall that Jost solution ${\bmi B}_{u}({\bmi x}^{\prime};\pm{\bmi K}_{T})$ 
from (\ref{Budef})  
was constructed to satisfy the {\em 
radiation condition\/} in the lower half space $z> z_L$, while 
Jost solution ${\bmi B}_{v}({\bmi x}^{\prime};\pm{\bmi K}_{T})$ from 
(\ref{Bvdef}) satisfies the {\em traction free condition\/} 
at the free surface $z=f$. At the poles, we assume ``one" of ``three" conditions is 
satisfied:
\begin{itemize}
\item the $P-$wave columns $(v,P,\pm{\bmi K}_{T}, z^{\prime})$ of 
${\bmi B}_{v}({\bmi x}^{\prime};\pm{\bmi K}_{T})$ satisfy the 
radiation condition
and the $P-$wave columns $(u,P,\pm{\bmi K}_{T}, z^{\prime})$ of 
${\bmi B}_{u}({\bmi x}^{\prime};\pm{\bmi K}_{T})$ satisfy the 
free surface condition.

\item the $SV-$wave columns $(v,S,\pm{\bmi K}_{T}, z^{\prime})$ of 
${\bmi B}_{v}({\bmi x}^{\prime};\pm{\bmi K}_{T})$ satisfy the 
radiation condition
and the $SV-$wave columns $(u,S,\pm{\bmi K}_{T}, z^{\prime})$ of 
${\bmi B}_{u}({\bmi x}^{\prime};\pm{\bmi K}_{T})$ satisfy the 
free surface condition.

\item the $SH-$wave columns $(v,H,\pm{\bmi K}_{T}, z^{\prime})$ of 
${\bmi B}_{v}({\bmi x}^{\prime};\pm{\bmi K}_{T})$ satisfy the 
radiation condition
and the $SH-$wave columns $(u,H,\pm{\bmi K}_{T}, z^{\prime})$ of 
${\bmi B}_{u}({\bmi x}^{\prime};\pm{\bmi K}_{T})$ satisfy the 
free surface condition.

\end{itemize}
By actually locating a surface wave pole, it can be verified that, either
\begin{itemize}
\item $(v,P,{\bmi K}_{T},z)$ is a linear multiple $(u,P,{\bmi K}_{T},z)$ 
\item or $(v,S,{\bmi K}_{T},z)$ is a linear multiple $(u,S,{\bmi K}_{T},z)$ 
\item or $(v,H,{\bmi K}_{T},z)$ is a linear multiple $(u,H,{\bmi K}_{T},z)$
\end{itemize}
The {\em degenerate\/} case where any two of the $P$, $S$, $H$ columns in 
${\bmi B}_{u}$ and ${\bmi  B}_{v}$ are scalar multiples, leads to a 
second order pole, and is not considered further, although this 
possibility cannot be ruled out {\em a priori\/}.


Following the method of {\em Kennett\/}~(1983) Ch.11, we employ Cramer's 
rule to write 
\begin{equation}  \label{InvCond}
\left\langle{\bmi  B}_{u}({\bmi x}^{\prime},-{\bmi K}_{T}),{\bmi 
B}_{v}({\bmi x}^{\prime},{\bmi K}_{T})\right\rangle^{-1}=\frac{
\mbox{Adj} \left\langle{\bmi  B}_{u}({\bmi x}^{\prime},-{\bmi K}_{T}),{\bmi 
B}_{v}({\bmi x}^{\prime},{\bmi K}_{T})\right\rangle }{\mbox{Det}\left\langle{\bmi  B}_{u}({\bmi x}^{\prime},-{\bmi K}_{T}),{\bmi 
B}_{v}({\bmi x}^{\prime},{\bmi K}_{T})\right\rangle}
\end{equation}
where 
\[
\mbox{Det}\left(\left\langle{\bmi  B}_{u}(z,-{\bmi K}_{T}),
{\bmi  B}_{v}({\bmi x}^{\prime},{\bmi K}_{T})\right\rangle\right)
\]
is the determinant and 
\[
\mbox{Adj}\left(\left\langle{\bmi  B}_{u}(z,-{\bmi K}_{T}),
{\bmi  B}_{v}({\bmi x}^{\prime},{\bmi K}_{T})\right\rangle\right)
\]
is the adjoint matrix (i.e. the transposed matrix of cofactors).


Let's assume $(x-x_S)>0$, and we have a P-wave pole at 
${\bmi K}_{T}={\bmi k}_{\mu}$ in the 
{\em first quadrant of the complex $k_x$ plane\/}, 
so that the P-wave columns satisfy 
\[
(u,P,-{\bmi k}_{\mu})=\gamma_{P} (v,P,-{\bmi k}_{\mu}),
\]
 for an arbitrary constant 
$\gamma_{P}$, then using similar arguments to those employed in {\em 
Kennett\/}(1983),  Appendix to chapter 11, we get:
\begin{eqnarray} \label{zeroinv}
\lefteqn{
\!\!\!\!\!\!
\left\langle{\left({u,P,-{\bmi K}_{T}}\right),\left({v,P,{\bmi K}_{T}}\right)}\right\rangle=
\left\langle{\left({u,P,-{\bmi K}_{T}}\right),\left({v,S,{\bmi K}_{T}}\right)}\right\rangle=
\left\langle{\left({u,P,-{\bmi K}_{T}}\right),\left({v,H,{\bmi 
K}_{T}}\right)}\right\rangle}\nonumber\\
&&
=\left\langle{\left({u,S,-{\bmi K}_{T}}\right),\left({v,P,{\bmi K}_{T}}\right)}\right\rangle=
\left\langle{\left({u,H,-{\bmi K}_{T}}\right),\left({v,P,{\bmi 
K}_{T}}\right)}\right\rangle=0
\end{eqnarray}
because both factors in $\langle \cdot, \cdot \rangle$ satisfy the same 
homogeneous boundary conditions (i.e. free surface or radiation). Also, 
the bilinear form  $\langle \cdot, \cdot \rangle$ defined on the Jost 
solutions is depth invariant, for the reasons explained in (\ref{invariant}), so 
that if a term like $\left\langle{\left({u,P,-{\bmi K}_{T}}\right),\left({v,P,{\bmi 
K}_{T}}\right)}\right\rangle$ vanishes at either the free surface $z=z_f$ 
or at infinity $z=\infty$ then $\left\langle{\left({u,P,-{\bmi K}_{T}}\right),\left({v,P,{\bmi 
K}_{T}}\right)}\right\rangle=0, \ z_f \le z \le \infty$ (which is why we 
drop the spatial dependence in Jost solutions like $\left({u,P,-{\bmi 
K}_{T}}\right)$ inside $\langle \cdot, \cdot \rangle$, 
when space on the printed page is at a premium). 


Employing (\ref{zeroinv}) in (\ref{InvCond}) and canceling common factors,
we get
\begin{equation} \label{zeroinv1}
\left\langle{\bmi  B}_{u}({\bmi x}^{\prime},-{\bmi K}_{T}),{\bmi 
B}_{v}({\bmi x}^{\prime},{\bmi K}_{T})\right\rangle^{-1}\sim 
\frac{
\left({\begin{array}{ccc}1&0&0\\0&0&0\\0&0&0\end{array}}\right)
}{\left\langle{\left({u,P,-{\bmi K}_{T}}\right),
\left({v,P,{\bmi K}_{T}}\right)}\right\rangle}
\ \mbox{as} \ {\bmi K}_{T}\rightarrow {\bmi k}_{\mu}
\end{equation}
where ${\bmi k}_{\mu}$ is a pole position in the complex wavenumber plane. 
Substituting (\ref{zeroinv1}) for  $\left\langle{\bmi  B}_{u}({\bmi x}^{\prime},-{\bmi K}_{T}),{\bmi 
B}_{v}({\bmi x}^{\prime},{\bmi K}_{T})\right\rangle^{-1}$ in the Green's tensor 
(\ref{GT0}), the singular contribution $\rf{\bmi G}_{P,\mu}$ at the $P-$wave pole
${\bmi K}_{T}= {\bmi k}_{\mu}$ may be written as
\begin{samepage}
\begin{eqnarray}\label{GTres}
\lefteqn{\rf{\bmi G}_{P,\mu}\!({\bmi x},{\bmi x}_{S};k_{y}) \sim 
\frac{1}{2 \pi}\int_{-\infty}^{\infty} {\rm d}k_x\, 
\exp\left[{i{k}_{x }({x-{x}_{S}})}\right]}\\
&&
\times\left\{
{\begin{array}{ll}
-\frac{\textstyle \left({ v,P,{\bmi K}_{T},z}\right) \otimes
\left({ u,P,-{\bmi K}_{T},z_S}\right) {\bmi  J}_{6}}
{\textstyle \left\langle{\left({u,P,-{\bmi K}_{T}}\right),
\left({v,P,{\bmi K}_{T}}\right)}\right\rangle}& 
z<z_{S} \\
-\frac{\textstyle \left({ u,P,{\bmi K}_{T},z}\right) \otimes
\left({ v,P,-{\bmi K}_{T},z_S}\right) {\bmi  J}_{6}}
{\textstyle \left\langle{\left({u,P,-{\bmi K}_{T}}\right),
\left({v,P,{\bmi K}_{T}}\right)}\right\rangle}& z>z_{S} 
\end{array} }\right. \quad \mbox{as} \quad {\bmi K}_{T}\rightarrow {\bmi k}_{\mu}
\nonumber
\end{eqnarray} 
\end{samepage}
but employing the linear dependence relation $(u,P,-{\bmi k}_{\mu})=
\gamma (v,P,-{\bmi k}_{\mu})$, in  (\ref{GTres})  and canceling common 
factors, we get
\begin{samepage}
\begin{eqnarray}\label{GTres1}
\lefteqn{\rf{\bmi G}_{P,\mu}\!({\bmi x},{\bmi x}_{S};k_{y}) \sim 
\frac{1}{2 \pi}\int_{-\infty}^{\infty} {\rm d}k_x\, 
\exp\left[{i{k}_{x }({x-{x}_{S}})}\right]}\\
&&
\times\left\{
{\begin{array}{ll}
-\frac{\textstyle \left({ v,P,{\bmi K}_{T},z}\right) \otimes
\left({ v,P,-{\bmi K}_{T},z_S}\right) {\bmi  J}_{6}}
{\textstyle \left\langle{\left({v,P,-{\bmi K}_{T}}\right),
\left({v,P,{\bmi K}_{T}}\right)}\right\rangle}& 
0 \le z,z_{S} \le \infty \\
-\frac{\textstyle \left({ u,P,{\bmi K}_{T},z}\right) \otimes
\left({ u,P,-{\bmi K}_{T},z_S}\right) {\bmi  J}_{6}}
{\textstyle \left\langle{\left({u,P,-{\bmi K}_{T}}\right),
\left({u,P,{\bmi K}_{T}}\right)}\right\rangle}& 
0 \le z,z_{S} \le \infty 
\end{array} }\right. \quad \mbox{as} \quad {\bmi K}_{T}\rightarrow {\bmi k}_{\mu}
\nonumber
\end{eqnarray} 
\end{samepage}
The surface wave is given 
by the $ 2 \pi i \times$ {residue of } (\ref{GTres1}) at $
{\bmi k}_{ \mu }=(k_{\mu, x}(k_{ y},\omega), k_{ y})$  in first quadrant
of the complex $k_x-$plane, so that
\begin{equation} 	\label{pole1}
	\rf{\bmi  G}_{P,\mu}\!({\bmi  x},{\bmi  x}_{S};k_{y})=
	-i {\frac{\exp\left[{i{k}_{ \mu,x }({x-{x}_{S}})}\right]
	\left({ u,P,{\bmi k}_{ \mu 	},{z}}\right)\otimes 
	\left({ u,P,-{\bmi k}_{ \mu },{z}_{S}}\right)
	{ \bmi  J}_{6} }{ {\frac{{\rm d}\,}{ {\rm d}\,{k}_{x}}} 
	{\left.{\left\langle{\left({ u,P,-{\bmi K}_{T}}\right) , 
	\left({u,P,{\bmi K}_{T}}\right)}\right\rangle}\right|}_{{\bmi K}_{T} =
	{\bmi k}_{ \mu }}}}  
\end{equation}

Similarly, assuming that $(x-x_S)<0$ we pick up $- 2 \pi i \times $ the 
residue  at the  pole  ${\bmi K}_{T}= - {\bmi k}_{\mu}=-(k_{\mu, x}(k_{ y},\omega), k_{ y})$, in the
third quadrant of the complex $k_x-$plane, so that
\begin{equation}	\label{pole2}
	\rf{\bmi  G}_{P,\mu}\!({\bmi  x},{\bmi  x}_{S};k_{y})=
	i\,
	\frac{ \exp \left[{i{k}_{ \mu,x }({x-{x}_{S}})}\right] \left({ u,P,{\bmi k}_{ \mu 
	},z}\right)\otimes \left({ u,P,-{\bmi k}_{ \mu },z_S}\right){ \bmi  J}_{6}}{ 
	{\frac{{\rm d}\,}{ {\rm d}\,{ k}_{x}}} {\left.{\left\langle{\left({ 
	u,P,-{\bmi K}_{T }}\right) ,\ \left({u,P,{\bmi K}_{T 
	}}\right)}\right\rangle}\right|}_{{\bmi K}_{T} =-{\bmi k}_{ \mu }}}  
\end{equation}
but it turns out that 
\[
{\frac{{\rm d}\,}{ {\rm d}\,{ k}_{x}}}{\left.\left\langle\left( 
	u,P,-{k}_{ x }\right) ,\ \left(u,P,{\bmi K}_{T} 
	\right)\right\rangle\right|}_{{\bmi K}_{T} =-{\bmi k}_{ \mu }} = - 
{\frac{{\rm d}\,}{ {\rm d}\,{ k}_{x}}}{\left.\left\langle\left( 
	u,P,-{\bmi K}_{T}\right) ,\ \left(u,P,{\bmi K}_{T}\right)
	\right\rangle\right|}_{{\bmi K}_{T} ={\bmi k}_{ \mu }}
\]
so that  (\ref{pole1}) and (\ref{pole2}) together imply that
\begin{equation} 	\label{polenu}
	\rf{\bmi  G}_{P,\mu}\!({\bmi  x},{\bmi  x}_{S};k_{y})=-i {\frac{
	\exp\left[{i{\bmi k}_{ \mu }|{x-{x}_{S}}|}\right] \left({ v,P,{\bmi k}_{ \mu 
	},z}\right)\otimes \left({ u,P,-{\bmi k}_{ \mu },z_S}\right){ \bmi  J}_{6}}{ 
	{\frac{{\rm d}\,}{ {\rm d}\,{ k}_{x}}} {\left.{\left\langle{\left({ 
	u,P,-{k}_{ x },\zeta}\right) ,\ \left({v,P,{k}_{ x 
	},\zeta}\right)}\right\rangle}\right|}_{{ k}_{ x} ={\bmi k}_{ \mu }}}}  
\end{equation}

We could calculate 
${\frac{{\rm d}\,}{ {\rm d}\,{ k}_{x}}}{\left.\left\langle\left( 
	u,P,-{k}_{ x },\zeta\right) ,\ \left(v,P,{k}_{ x 
	},\zeta\right)\right\rangle\right|}_{{ k}_{ x} ={\bmi k}_{ \mu }}$ 
in (\ref{polenu}) using a finite difference approximation to the 
derivative, but the calculation is rather susceptible to roundoff. 
Instead, we develop an analytic expression, that is 
the cartesian  analogue  of the result in {\em 
Kennett\/}~(1983),  Appendix to chapter 11.

\subsection{Calculating the residue at a surface wave pole}
Since
\begin{equation} 	\label{m1}
	{ \partial }_{z}\left(u,P,\pm{\bmi K}_{T},z\right)=
	{\rf{\bmi A}}(z,\pm{\bmi K}_{T})\left(u,P,\pm{\bmi K}_{T},z\right)
\end{equation}
where  ${\rf{\bmi A}}(z,{\bmi K}_{T})$ is the Fourier domain system matrix
\begin{eqnarray} \label{Akzdef}
\lefteqn{
{\rf{\bmi A}}(z,{\bmi K}_{T})
=}\\
&&\!\!\!\!\!\!\!\!\!\!\!\!
\left[{\begin{array}{cccccc}
0&- i k_x {\rf \gamma}(z) &  -i {  k}_{  y} {\rf \gamma}(z) &  {\rf a}(z)&0&0\\ 
-i k_x&0&0&0&{\rf b}(z)&0\\ -i{k}_{y}&0&0&0&0&{\rf b}(z)\\ - {\rf \rho}(z) 
{\omega }^{  2}&  0&0&0&-i k_x&-i{k}_{y}\\ 0&\!\!\!\!\!\!
k^{2}_x{\rf \zeta}(z)  + {\rf \mu}(z) {  k}_{  y}^{  2}- {\rf \rho}(z) {\omega 
}^{  2} &  
{k}_{x}{k}_{y} ({\rf \chi}(z)+ {\rf \mu}(z))&  -i {k}_{x} {\rf \gamma}(z) &  0&0\\ 
0&{k}_{x}{k}_{y} ({\rf \chi}(z)+ {\rf \mu}(z)) & \!\!\! \!\!\!
{\rf \zeta }(z){  k}_{  y}^{  2}  + {  k}_{  x}^{  2} {\rf \mu}(z) 
- {\rf \rho}(z) {\omega }^{  2}   &  -i{k}_{y} {\rf \gamma}(z) &  
0&0\end{array}}\right]\nonumber
\end{eqnarray}
it follows that
\begin{eqnarray} 	\label{m2}
{\partial }_{z}{\frac{{\rm d}\,}{{\rm 
d}\,{ k}_{x}}}\left({u,P,\pm{\bmi K}_{T},z}\right)&=&
\rf{\bmi  A}(z,\pm{\bmi K}_{T}){\frac{{\rm d}\,}{{\rm 
d}\,{ k}_{x}}}\left({u,P,\pm{\bmi K}_{T},z}\right)\nonumber\\
&&
+
{\frac{{\rm d}\,}{{\rm 
d}\,{ k}_{x}}}\rf{\bmi  A}(z,\pm{\bmi K}_{T})\left({u,P,\pm{\bmi 
K}_{T},z}\right)
\end{eqnarray}
Using  (\ref{m1}) and (\ref{m2}) gives
\begin{eqnarray} \label{m4}
\lefteqn{{ \partial 
}_{z}\left\langle{{\frac{{\rm d}\,}{{\rm d}\,{ k}_{x}}}
\left({u,P,-{\bmi K}_{T},z}\right), 
\left({u,P,{\bmi K}_{T},z}\right)}\right\rangle}\nonumber\\ 
&&=\left\langle{{\rm \partial 
}_{z}{\frac{{\rm d}\,}{{\rm d}\,{ k}_{x}}}\left({u,P,-{\bmi K}_{T},z}\right), 
\left({u,P,{\bmi K}_{T},z}\right)}\right\rangle+\left\langle 
{{\frac{{\rm d}\,}{{\rm d}\,{ k}_{x}}}\left({u,P,-{\bmi K}_{T},z}\right),{\rm \partial 
}_{z}\left({u,P,{\bmi K}_{T},z}\right)}\right\rangle\nonumber\\
&&=\left\langle{\rf{\bmi A}(z,-{\bmi K}_{T}){\frac{{\rm d}\,}{{\rm d}\,{ k}_{x}}}\left({u,P,-{\bmi K}_{T},z}\right),
\left({u,P,{\bmi K}_{T},z}\right)}\right\rangle\nonumber\\
&&\qquad+\left\langle{
{\frac{{\rm d}\,}{{\rm 
d}\,{ k}_{x}}}\rf{\bmi  A}(z,-{\bmi K}_{T})
\left({u,P,-{\bmi K}_{T},z}\right),\left({u,P,{\bmi K}_{T},z}\right)}\right\rangle\nonumber\\
&&\qquad +
\left\langle{{\frac{{\rm d}\,}{{\rm d}\,{ k}_{x}}}
\left({u,P,-{\bmi K}_{T},z}\right),\rf{\bmi A}(z,{\bmi K}_{T})
\left({u,P,{\bmi K}_{T},z}\right)}\right\rangle\nonumber\\
&&=-\left\langle{
\left({u,P,-{\bmi K}_{T},z}\right),{\frac{{\rm d}\,}{{\rm 
d}\,{ k}_{x}}}\rf{\bmi  A}(z,{\bmi K}_{T})\left({u,P,{\bmi K}_{T},z}\right)}\right\rangle
\end{eqnarray} 
where we have used the Hamiltonian symmetries
\[
\left\langle{\rf{\bmi  A}(z,-{\bmi K}_{T}) \cdot   , \cdot 
}\right\rangle=-\left\langle{ \cdot   ,\ \rf{\bmi  A}(z,{\bmi K}_{T}) 
\cdot }\right\rangle
\]
and
\[
\langle{{\frac{{\rm d}\,}{{\rm d}\,{ k}_{x}}}
\rf{\bmi  A}(z,-{\bmi K}_{T}) \cdot   
, \cdot }\rangle=-\langle{ \cdot   ,\ 
{\frac{{\rm d}\,}{{\rm 
d}\,{ k}_{x}}}\rf{\bmi  A}(z,{\bmi K}_{T}) \cdot }\rangle
\] 
Upon integrating (\ref{m4}) over the depth interval $z\le \zeta \le 
\infty$, 
and using the {\em radiation condition\/} at $\zeta =\infty$, we get
\begin{eqnarray}	\label{mint1}
\lefteqn{
{\left.{\left\langle{{\frac{{\rm d}\,}{{\rm d}\,{ k}_{ x}}}\left({ 
	u,P,-{\bmi K}_{T},\zeta }\right) ,\left({u,P,{\bmi K}_{T},\zeta 
	}\right)}\right\rangle}\right|}_{ \zeta   =z} =}\nonumber\\
&&\int_{ 
	z}^{\infty } {\rm d} \zeta \left\langle{\left({ u,P,-{\bmi K}_{T},\zeta 
	}\right) ,{\frac{{\rm d}\,}{{\rm 
d}\,{ k}_{x}}}\rf{\bmi  A}(\zeta  ,{\bmi K}_{T})\left({u,P,{\bmi K}_{T},\zeta 
	}\right)}\right\rangle
\end{eqnarray}




Similarly, integrating 
(\ref{m4}) over the depth interval $0\le \zeta \le z$, and remembering that
$\left({ u,P,\pm{\bmi K}_{T},\zeta=0 }\right)$ satisfy the traction free 
condition at $\zeta =0$ in the limit as $ \pm{\bmi K}_{T} 
\rightarrow \pm{\bmi k}_{\mu}$ (which is the only case of interest),
we get 
\begin{eqnarray} \label{mint2}
\lefteqn{{\left.{\left\langle{\left({ u,P,-{\bmi K}_{T},\zeta }\right) , 
{\frac{{\rm d}\,}{{\rm d}\,{ 
k}_{ x}}}\left({u,P,{\bmi K}_{T},\zeta }\right)}\right\rangle}\right|}_{ \zeta 
=z} =}\nonumber\\
&&\int_{ 0}^{z } {\rm d} \zeta \left\langle{\left({ u,P,-{\bmi K}_{T},\zeta 
}\right) ,{\frac{{\rm d}\,}{{\rm 
d}\,{ k}_{x}}}\rf{\bmi  A}(\zeta ,{\bmi K}_{T})\left({u,P,{\bmi K}_{T},\zeta 
}\right)}\right\rangle 
\end{eqnarray} 
Now, (\ref{mint1}) and 
(\ref{mint2}) together imply that 
\begin{eqnarray}	\label{deriv}
\lefteqn{{\frac{{\rm d}\,}{{\rm d}\,{ k}_{ 
	x}}}\left\langle{\left({ u,P,-{\bmi K}_{T},{ z}}\right)
	,\left({u,P,{\bmi K}_{T},{ z}}\right)}\right\rangle =}\nonumber\\
&&\int_{ 
	0}^{\infty } {\rm d} \zeta \left\langle{\left({ u,P,-{\bmi K}_{T},\zeta 
	}\right) ,{\frac{{\rm d}\,}{{\rm 
d}\,{ k}_{x}}}\rf{\bmi  A}(\zeta  ,{\bmi K}_{T})\left({u,P,{\bmi K}_{T},\zeta 
	}\right)}\right\rangle
\end{eqnarray}
for an arbitrary depth point $z$.

\subsection{The surface wave contributions to the Green's tensor}

 Finally, using (\ref{deriv}) in (\ref{polenu}), and summing over the
 $N_{P}(k_y, \omega )$ $P-$wave poles in the first quadrant of the $k_x$ plane 
at fixed $(k_y, \omega)$, 
we get the surface wave contributions to the Green's tensor as
\begin{eqnarray}	\label{surfGp}
\lefteqn{
\rf{\bmi  G}_{P}({\bmi  x},{\bmi  x}_{S};k_y)=-i 
\sum\limits_{\mu = 1}^{  N_{P}( \omega)} 
{\frac{\exp\left[{i{ k}_{ \mu,x }\left|{x-{x}_{S}}\right|}\right]
	\left({  u,P,{\bmi k}_{ \mu },z}\right)\otimes \left({  
	u,P,-{\bmi k}_{ \mu },z_S}\right){ \bmi  J}_{6}}{\int_{  
	0}^{\infty }  {\rm d} \zeta {\left.{\left\langle{\left({  
	u,P,-{\bmi K}_{T},\zeta}\right)  ,{\frac{{\rm d}\,}{{\rm 
d}\,{ k}_{x}}}\rf{\bmi  A}(\zeta,k_x)
	\left({u,P,{\bmi K}_{T},\zeta}\right)}\right\rangle}\right|}_{{ 
	 \bmi K}_{T}  ={\bmi k}_{ \mu }}}}  
}\nonumber\\
&& 
\end{eqnarray}

The surface wave contributions to  $\rf{\bmi  G}_{S}$ and  $\rf{\bmi  G}_{H}$
 that result from the 
\[
(u,S,-{\bmi k}_{\mu})=\gamma_{S} (v,S,-{\bmi k}_{\mu}),
\]
and
\[
(u,H,-{\bmi k}_{\mu})=\gamma_{H} (v,H,-{\bmi k}_{\mu}),
\]
cases respectively,  are treated  with the same method that we have just outlined for 
$\rf{\bmi  G}_{P}$ in (\ref{surfGp}), although the pole positions ${\bmi 
k}_{\mu}$ will vary in each case.

To summarize, the surface wave Green's tensor may be written as 
\begin{eqnarray}	\label{surfG}
\lefteqn{
\rf{\bmi  G}({\bmi  x},{\bmi  x}_{S};k_y)=-i \times
}\\
&& \sum\limits_{c \in \{ P,S,H \}}
\sum\limits_{\mu_c = 1}^{  N_{c}( \omega)}  
{\frac{\exp\left[{i{ k}_{ \mu_c,x }\left|{x-{x}_{S}}\right|}\right]
	\left({  u,P,{\bmi k}_{ \mu_c },z}\right)\otimes \left({  
	u,P,-{\bmi k}_{ \mu_c },z_S}\right){ \bmi  J}_{6}}{\int_{  
	0}^{\infty }  {\rm d} \zeta {\left.{\left\langle{\left({  
	u,P,-{\bmi K}_{T},\zeta}\right)  ,{\frac{{\rm d}\,}{{\rm 
d}\,{ k}_{x}}}\rf{\bmi  A}(\zeta,k_x)
	\left({u,P,{\bmi K}_{T},\zeta}\right)}\right\rangle}\right|}_{{ 
	 \bmi K}_{T}  ={\bmi k}_{ \mu_c }}}}\nonumber 
\end{eqnarray}

We recall that the specific formula chosen in (\ref{GT0}) depends on 
whether $(z-z_S) \gtlt 0$ ( the {\em up/down dichotomy\/}), 
while the specific formula for the extracted 
surface wave components (\ref{surfG}) depends on
whether $(x-x_S) \gtlt 0$ (the {\em left/right dichotomy\/}). 
Note that the complex plane portion wavenumber contour  in 
Fig.\ref{contour}
is subject to both the  {\em up/down dichotomy\/} and the {\em left/right 
dichotomy\/} which may be handled by an {\em alternating direction, Riccati 
sweep\/}, i.e. {\em each offset of the left/right sweep has an embedded 
up/down sweep\/}. The possibility of explicitly calculating modal/body wave 
coupling may be of use in certain studies.

Unfortunately, the stumbling block is the actual calculation of the 
surface wave pole positions. By employing the complex plane contour
of Fig.\ref{contour}, the number of plane waves needed to synthesize the 
body wave phases is greatly reduced, but the number of poles is very 
large indeed, and it is a major computational problem to locate them to 
the required accuracy.






